\documentclass[12pt]{article}

\usepackage[paperwidth=200mm,paperheight=112.5mm,margin=11mm,top=10mm]{geometry}
\usepackage{amsmath,amssymb,booktabs,array}
\usepackage[T1]{fontenc}
\usepackage{lmodern}

\pagestyle{empty}
\setlength{\parindent}{0pt}
\setlength{\abovedisplayskip}{9pt}
\setlength{\belowdisplayskip}{9pt}

\newcommand{\R}{\mathbb{R}}
\newcommand{\T}{\mathsf{T}}
\DeclareMathOperator{\col}{col}
\DeclareMathOperator{\range}{range}
\newcommand{\fbase}{f_{\mathrm{base}}}
\newcommand{\faug}{f_{\mathrm{aug}}}
\newcommand{\Faug}{F_{\mathrm{aug}}}
\newcommand{\ANN}{\mathrm{ANN}}
\newcommand{\xio}{x_k^{\mathrm{IO}}}
\newcommand{\thaux}{\theta_{\mathrm{aux}}}
\newcommand{\panel}[1]{\subsection*{\Large #1}\vspace{1mm}}

\begin{document}

%======================================================================
%======================================================================
\panel{Gy\"or\"ok: orthogonal by construction}

His baseline is exactly linear in its parameters, so stacking over the data
set gives $\Phi\in\R^{Nn_y\times n_{\theta_b}}$ with $\widehat Y=\Phi\theta_b
+F^{\ANN}_{\theta_a}$ (Eq.~5), where $\xio$ is the lagged input--output
regressor. He then defines

\begin{equation*}
\thaux=\left(\Phi^{\T}\Phi\right)^{-1}\Phi^{\T}F^{\ANN}_{\theta_a},
\tag{Eq.~8}
\end{equation*}

\vspace{-2mm}
\begin{equation*}
\widetilde F^{\ANN}_{\theta_a}
=F^{\ANN}_{\theta_a}-\Phi\thaux
=\left[I-\Phi\left(\Phi^{\T}\Phi\right)^{-1}\Phi^{\T}\right]F^{\ANN}_{\theta_a},
\tag{Eq.~10}
\end{equation*}

\vspace{-2mm}
\begin{equation*}
\Phi^{\T}\widetilde F^{\ANN}_{\theta_a}=0 .
\tag{Lemma~3, Eq.~11}
\end{equation*}

$\thaux\in\R^{n_{\theta_b}}$: the least-squares fit of the ANN output onto
$\Phi$.

\vspace{1mm}
$\Phi$ is built once from the training regressors and never changes, since
$\phi$ does not depend on $\theta_b$. $\thaux$ is recomputed every time the
loss is evaluated.

\newpage
%======================================================================
\panel{The baseline is not linear in its parameters}

Gy\"or\"ok (2025) expands the state transition around a linearization point
$\bar\theta$:

\begin{equation*}
f_\theta(x_k,u_k)\approx f_{\bar\theta}(x_k,u_k)
+\Phi_{\bar\theta}(x_k,u_k)\left(\theta-\bar\theta\right),
\qquad
\Phi_{\bar\theta}(x_k,u_k)
=\left.\frac{\partial f_\theta(x_k,u_k)}{\partial\theta}\right|_{\theta=\bar\theta}
\in\R^{n_x\times n_\theta}.
\tag{Eq.~15}
\end{equation*}

Applied to the gantry transition, with the ten physical free parameters $v$ and
the current estimate $\bar v$:

\begin{equation*}
\fbase(v,x_{b,k},u_k)\approx\fbase(\bar v,x_{b,k},u_k)
+\Phi_{\bar v}(x_{b,k},u_k)\left(v-\bar v\right),
\qquad
\Phi_{\bar v}(x_{b,k},u_k)
=\left.\frac{\partial \fbase}{\partial v}\right|_{v=\bar v}
\in\R^{6\times 10}.
\end{equation*}

Gy\"or\"ok (2025) uses $\Phi_{\bar\theta}$ inside the $\beta$-weighted penalty
(Eq.~13); here it is the span in the projection of Eq.~8--11.

\newpage
%======================================================================
\panel{The reference set that $J$ is built on}

Fixed one-step evaluation points from the training records, not a simulated
trajectory. Velocities are unmeasured, since $C_y=\begin{bmatrix}P^{\T}&0
\end{bmatrix}$, so they are reconstructed:

\begin{equation*}
q_k=P^{-\T}y_k,
\qquad
\dot q_k=\frac{-q_{k+2}+8q_{k+1}-8q_{k-1}+q_{k-2}}{12T_s},
\qquad
x_{a,k}=0,
\qquad
u_k\ \text{recorded}.
\end{equation*}

Two samples are trimmed at each record end so the stencil never crosses a
boundary. Stacking sample-major over the routed physical rows,

\begin{equation*}
J=\col_k J_{b,k}(\bar v)\in\R^{Nn_r\times 10},
\qquad
\Faug(\theta_a)=\col_k \faug^{b}(\theta_a,x_{b,k},0,u_k)\in\R^{Nn_r},
\qquad
N=671\,944 .
\end{equation*}

The tuples stay fixed for the whole run. The set has numerical rank ten, so all
ten parameter directions are excited.

\newpage
%======================================================================
\panel{The same construction, with $J$ in place of $\Phi$}

\begin{equation*}
J=U\Sigma V^{\T},
\end{equation*}

\vspace{-2mm}
\begin{equation*}
\thaux(\theta_a)=J^{\dagger}\Faug(\theta_a)=V\Sigma^{-1}U^{\T}\Faug(\theta_a)
\in\R^{10},
\tag{cf.\ Eq.~8}
\end{equation*}

\vspace{-2mm}
\begin{equation*}
\widetilde\Faug=\Faug-J\thaux,
\tag{cf.\ Eq.~10}
\end{equation*}

\vspace{-2mm}
\begin{equation*}
J^{\T}\widetilde\Faug=0 .
\tag{cf.\ Eq.~11}
\end{equation*}

The Gram matrix is never formed; the solve uses the economy SVD in double
precision. The inner product is the Euclidean one in the model's normalized
state coordinates.

\newpage
%======================================================================
\panel{The correction inside the rollout}

\begin{equation*}
\begin{aligned}
x_{b,k+1}&=\fbase(v,x_{b,k},u_k)+\faug^{b}(\theta_a,\widehat x_k,u_k)
-J_b(\bar v,x_{b,k},u_k)\,\thaux,\\[1mm]
x_{a,k+1}&=\faug^{a}(\theta_a,\widehat x_k,u_k).
\end{aligned}
\end{equation*}

Gy\"or\"ok evaluates his stored coefficient against the regressor at the new
point, $\widehat y_k=\phi(\xio)\widehat\theta_b+f^{\ANN}(\xio)
-\phi(\xio)\widehat\thaux$ (Eq.~13). Here $\thaux$ is fitted on the reference
tuples but applied at the rollout state, with the epoch expansion point
$\bar v$ frozen in both places.

\vspace{1mm}
Lifecycle: the reference tuples are fixed for the run; $\bar v$ is detached and
$J$ rebuilt at each epoch boundary; $\thaux$ is recomputed at every objective
evaluation and never reused across optimizer updates.

\newpage
%======================================================================
\panel{What this establishes}

\begin{equation*}
\text{established:}\qquad
J^{\T}\widetilde\Faug=0
\quad\text{on the fixed reference tuples, at }x_a=0,\ \text{at }\bar v .
\end{equation*}

\begin{equation*}
\text{not established:}\qquad
x_{a,k}\ \longrightarrow\ \faug^{b}
\ \longrightarrow\ x_{b,k+1}
\ \longrightarrow\ \cdots\ \longrightarrow\ y_{k+n}
\end{equation*}

The condition is one-step, stacked, in normalized routed physical-state
coordinates, against the local ten-parameter tangent at $\bar v$. It is not
pointwise, not global in $v$, not valid at nonzero $x_a$, and says nothing
about the accumulated input--output Jacobian after propagation and feedback.
Gy\"or\"ok's parameter recovery (Theorem~7) needs an exact $\Phi$ and
Condition~4, $\sum_i\phi^{\T}(\xio)\delta(\xio)=0$, neither of which transfers
automatically.

\end{document}
